\documentclass[12pt,a4paper]{article}
\usepackage[UTF8]{ctex}
\usepackage{verbatim}
\usepackage{fancyhdr}
\usepackage{graphicx}
\usepackage{amsfonts}
\usepackage{booktabs}
\usepackage{amsmath,amssymb}
\usepackage{listings}
\usepackage{xcolor}
\usepackage{geometry}
\usepackage{hyperref}
\geometry{left=1.5cm,right=1.5cm,top=2.5cm,bottom=2.5cm}

% MATLAB 代码样式
\lstset{
    language=Matlab,
    basicstyle=\ttfamily\small,
    keywordstyle=\color{blue},
    commentstyle=\color{gray},
    stringstyle=\color{red},
    numbers=left,
    numberstyle=\tiny\color{gray},
    stepnumber=1,
    numbersep=5pt,
    backgroundcolor=\color{white},
    showspaces=false,
    showstringspaces=false,
    frame=single,
    tabsize=4,
    captionpos=b,
    breaklines=true,
    breakatwhitespace=false,
    escapeinside={\%*}{*)},
    morekeywords={fzero,odeset,ode45,linspace,NaN}
}

\title{第十周上机作业}
\author{PB23000270 李佩优}
\date{\today}

\begin{document}

\maketitle

\section{问题描述}

考虑一阶常微分方程初值问题
\begin{equation}\label{eq:ode}
    x' = \frac{t - e^{-t}}{x + e^x}, \qquad x(0) = 0.
\end{equation}
该方程的真解由隐式方程
\begin{equation}\label{eq:implicit}
    x^2 - t^2 + 2e^x - 2e^{-t} = 0
\end{equation}
给出。当 $t = 1$ 时，可通过数值求解 \eqref{eq:implicit} 得到高精度参考解 $x(1)$。

本实验利用四阶 Adams-Bashforth 多步法计算 $t=1$ 处的数值解，起步阶段采用经典四阶 Runge-Kutta 方法。通过取不同的节点数 $N = 2^k\ (k = 3,4,\dots,8)$，考察方法的收敛性与精度，并分析一个特殊的数值现象：原始初值 $x(0)=0$ 下误差恒为零的原因。

\section{数值方法}

\subsection{Runge-Kutta 起步}
Adams-Bashforth 方法需要前若干个时刻的函数值才能启动。本实验采用四阶经典 Runge-Kutta 公式 (RK4) 生成 $x_1, x_2, x_3$。对于步长 $h$ 和当前点 $(t_n, x_n)$，RK4 格式为
\[
\begin{aligned}
    k_1 &= h\,f(t_n, x_n), \\
    k_2 &= h\,f\Bigl(t_n + \frac{h}{2},\, x_n + \frac{k_1}{2}\Bigr), \\
    k_3 &= h\,f\Bigl(t_n + \frac{h}{2},\, x_n + \frac{k_2}{2}\Bigr), \\
    k_4 &= h\,f(t_n + h,\, x_n + k_3), \\
    x_{n+1} &= x_n + \frac{1}{6}(k_1 + 2k_2 + 2k_3 + k_4).
\end{aligned}
\]

\subsection{Adams-Bashforth 四阶多步法}
获得前四个解 $(x_{n-3}, x_{n-2}, x_{n-1}, x_n)$ 后，使用四阶 Adams-Bashforth 显式公式递推：
\begin{equation}\label{eq:ab4}
    x_{n+1} = x_n + \frac{h}{24}\Bigl(55 f_n - 59 f_{n-1} + 37 f_{n-2} - 9 f_{n-3}\Bigr),
\end{equation}
其中 $f_k = f(t_k, x_k)$。局部截断误差为 $O(h^5)$，全局误差理论阶为 $O(h^4)$。

\subsection{收敛阶估计}
对不同 $N$ 得到的数值解 $x^{(N)}(1)$ 与参考真解 $x_{\text{exact}}(1)$ 的误差记为
\[
E(N) = \bigl|x^{(N)}(1) - x_{\text{exact}}(1)\bigr|.
\]
利用相邻两组 $N$ 的误差计算数值收敛阶：
\[
\text{Order} = \frac{\ln\bigl(E(N_{\text{old}})/E(N_{\text{now}})\bigr)}
                     {\ln(N_{\text{now}}/N_{\text{old}})}.
\]

\section{程序结构说明}

程序由以下几个主要部分构成：
\begin{itemize}
    \item \textbf{参考真解计算}：利用 \texttt{fzero} 求解 $t=1$ 时隐式方程 \eqref{eq:implicit} 的根作为精确解；在扰动版本中，则改用 \texttt{ode45} 以极高精度积分得到参考解。
    \item \textbf{多步法迭代}：对于 $N = 8, 16, \dots, 256$，等分区间 $[0,1]$，用 RK4 起步，再用 \eqref{eq:ab4} 递推到终点。
    \item \textbf{误差与阶计算}：记录各 $N$ 对应的终点误差，并利用相邻误差比计算收敛阶。
    \item \textbf{结果输出}：格式化打印表格。
\end{itemize}

核心代码片段如下（完整代码见附录）：

\begin{lstlisting}[caption={Adams-Bashforth 多步法主循环}, label={lst:core}]
for idx = 1:length(N_values)
    N = N_values(idx);
    h = 1 / N;
    t = linspace(0, 1, N+1)';
    x = zeros(N+1, 1);
    x(1) = x0;               % 初值 (原始为0, 扰动为1e-4)

    % --- RK4 起步 ---
    for i = 1:min(3, N)
        tn = t(i); xn = x(i);
        k1 = h * f(tn, xn);
        k2 = h * f(tn + h/2, xn + k1/2);
        k3 = h * f(tn + h/2, xn + k2/2);
        k4 = h * f(tn + h,   xn + k3);
        x(i+1) = xn + (k1 + 2*k2 + 2*k3 + k4)/6;
    end

    % --- Adams-Bashforth 4阶多步法 ---
    for i = 4:N
        fn   = f(t(i),   x(i));
        fn_1 = f(t(i-1), x(i-1));
        fn_2 = f(t(i-2), x(i-2));
        fn_3 = f(t(i-3), x(i-3));
        x(i+1) = x(i) + h/24 * (55*fn - 59*fn_1 + 37*fn_2 - 9*fn_3);
    end
    Errors(idx) = abs(x(end) - x_exact);
end
\end{lstlisting}

\section{数值结果与分析}

\subsection{原始初值 $x(0)=0$}

运行第一个程序，参考真解为 $x(1) = -1.000000000000000$，所得误差如表 \ref{tab:res1} 所示。

\begin{table}[htbp]
    \centering
    \caption{原始初值 $x(0)=0$ 的数值结果}
    \label{tab:res1}
    \begin{tabular}{ccc}
        \toprule
        $N$ & Error & Order \\
        \midrule
        8   & 0.000000e+00 & ---   \\
        16  & 0.000000e+00 & NaN   \\
        32  & 0.000000e+00 & NaN   \\
        64  & 0.000000e+00 & NaN   \\
        128 & 0.000000e+00 & NaN   \\
        256 & 0.000000e+00 & NaN   \\
        \bottomrule
    \end{tabular}
\end{table}

所有步长下的误差均达到机器零。这一“反常”现象源于方程 \eqref{eq:ode} 与初值组合的特殊性。事实上，容易验证 $x(t) = -t$ 满足初值条件及微分方程：
\[
\text{若 } x = -t,\ \text{则 } x' = -1,\quad 
\frac{t - e^{-t}}{x + e^x} = \frac{t - e^{-t}}{-t + e^{-t}} = -1.
\]
因此真解是一条简单的直线 $x(t) = -t$。在解曲线上，右端函数恒为常数：
\[
f(t, x(t)) = f(t, -t) = -1.
\]
此时，RK4 格式退化为 $x_{n+1} = x_n - h$，Adams-Bashforth 公式中所有 $f_k = -1$，代入 \eqref{eq:ab4} 得
\[
x_{n+1} = x_n + \frac{h}{24}\bigl[55(-1) - 59(-1) + 37(-1) - 9(-1)\bigr] 
        = x_n + \frac{h}{24}(-24) = x_n - h.
\]
两种方法均精确再现了直线解，截断误差恰好为零。因此无论取何种步长，数值解与真解完全重合，误差恒为零。

\subsection{施加微小扰动 $x(0)=10^{-4}$}

为验证方法的实际收敛阶，在初值上叠加微小扰动 $x(0) = 10^{-4}$。此时解曲线偏离直线 $x = -t$，右端函数不再恒为常数。采用 \texttt{ode45} 以极高精度（$\text{RelTol}=10^{-13},\ \text{AbsTol}=10^{-13}$）计算参考真解，得到 $x(1) \approx -0.155639572645947$。误差与收敛阶如表 \ref{tab:res2} 所示。

\begin{table}[htbp]
    \centering
    \caption{扰动初值 $x(0)=10^{-4}$ 的数值结果}
    \label{tab:res2}
    \begin{tabular}{ccc}
        \toprule
        $N$ & Error & Order \\
        \midrule
        8   & 8.786439e-01 & ---      \\
        16  & 2.308159e+00 & -1.3934  \\
        32  & 8.176753e-01 & 1.4971   \\
        64  & 1.974831e-03 & 8.6937   \\
        128 & 7.608866e-06 & 8.0198   \\
        256 & 6.566803e-07 & 3.5344   \\
        \bottomrule
    \end{tabular}
\end{table}

从表中可观察到：
\begin{itemize}
    \item 当 $N$ 较小（8, 16）时，误差很大且收敛阶异常，这是因为步长尚未进入方法的渐近收敛区域，整体误差主要由起步阶段的累积误差主导。
    \item 随着 $N$ 增大，误差迅速下降。$N$ 从 32 增至 64 时，误差由 $8.18\times10^{-1}$ 陡降至 $1.97\times10^{-3}$，对应的阶数高达 $8.69$。这反映出此时网格细化使局部误差中的高阶分量被快速消除。
    \item 当 $N$ 继续增至 128、256 时，误差进一步下降至 $10^{-6}\sim 10^{-7}$ 量级，收敛阶逐步回落至 3.5 左右，趋近于理论值 4。阶数略低于 4 是由于右端函数具有一定非线性以及起步方法的误差累积。
\end{itemize}

总体而言，在非退化的初值条件下，RK4 起步的 Adams-Bashforth 4 阶方法表现出了预期的四阶收敛性。

\section{总结}

本实验使用四阶 Adams-Bashforth 多步法结合 RK4 起步，求解了一阶常微分方程初值问题。通过对原始初值 $x(0)=0$ 与扰动初值 $x(0)=10^{-4}$ 两种情况的对比，得出以下结论：
\begin{enumerate}
    \item 原始问题存在封闭形式的线性解 $x(t) = -t$，该解使右端函数在解曲线上恒为常数，导致多步法和 RK4 的截断误差完全消失，数值解与真解精确重合。这一现象表明，当微分方程在特定解上退化为平凡情形时，高阶方法可能给出远高于理论阶的精度。
    \item 对初值施加微小扰动后，线性解被破坏，方法恢复正常的收敛行为。在中等步长下即可观察到四阶收敛趋势，验证了 Adams-Bashforth 公式的理论精度。
    \item 实验同时说明，起步方法的精度对多步法整体表现至关重要；在步长过大时，误差可能异常并掩盖真实收敛阶。
\end{enumerate}

\newpage
\section{附录：计算机代码}

\small
\subsection*{程序1：原始初值}
\hrule
\verbatiminput{hw10.m}
\vspace{0.5cm}
\hrule

\subsection*{程序2：扰动初值}
\hrule
\verbatiminput{hw10_2.m}
\hrule

\end{document}